\documentclass[conference]{IEEEtran}
\IEEEoverridecommandlockouts

\usepackage{cite}
\usepackage{amsmath,amssymb,amsfonts,amsthm}
\usepackage{algorithmic}
\usepackage{algorithm}
\usepackage{graphicx}
\usepackage{textcomp}
\usepackage{xcolor}
\usepackage{booktabs}
\usepackage{multirow}
\usepackage{url}
\usepackage{balance}

% Theorem environments
\newtheorem{theorem}{Theorem}
\newtheorem{definition}{Definition}

\begin{document}

\title{Multi-Resource Temporal Interval Scheduling and Concurrency-Preserving Web Architectures for Complex Clinical Regimens}

\author{
\IEEEauthorblockN{Ayush Malviya}
\IEEEauthorblockA{\textit{Dept. of Computer Science and Engineering} \\
\textit{Faculty of Engineering and Technology}\\
AyurSutra Research Initiative\\
Email: ayush@example.com}
\and
\IEEEauthorblockN{[Co-Author Name]}
\IEEEauthorblockA{\textit{Dept. of Computer Science and Engineering} \\
\textit{Faculty of Engineering and Technology}\\
AyurSutra Research Initiative\\
Email: coauthor@example.com}
\and
\IEEEauthorblockN{[Faculty Advisor / Guide Name]}
\IEEEauthorblockA{\textit{Dept. of Computer Science and Engineering} \\
\textit{Faculty of Engineering and Technology}\\
AyurSutra Research Initiative\\
Email: advisor@example.com}
}

\maketitle

\begin{abstract}
Standard outpatient scheduling systems rely on single-server queuing formulations that model clinical appointments as independent, scalar interactions between a single clinician and a patient within an unconstrained office space. However, specialized clinical day-care procedures---such as surgical day-theatres, multi-day restorative regimens, hemodialysis, and infusion therapeutics---mandate the simultaneous synchronization of three coupled, heterogeneous resource pools: a specialized physical suite, a certified clinical practitioner matching procedural qualifications, and the patient. In addition, treatments span 30 to 90 minutes and are governed by strict multi-day sequential protocols. When deployed in stateless, distributed web environments, conventional booking engines suffer from catastrophic scheduling collisions, room contention, and Time-Of-Check to Time-Of-Use (TOCTOU) race conditions.

This paper presents the formalization, algorithmic design, and empirical evaluation of \textbf{AyurSutra}, a multi-resource constraint satisfaction scheduling engine and high-concurrency web architecture developed under the Design Science Research Methodology (DSRM). We model clinical scheduling as an interval graph overlap problem governed by Allen's temporal interval algebra and formulate a discrete constraint satisfaction algorithm that evaluates multi-entity feasibility across dynamic operating shifts. To eliminate TOCTOU vulnerabilities without the latency overhead of coarse-grained database locks, the architecture introduces a write-time relational boundary revalidation protocol that provides atomic collision mitigation at the storage interface. Built on a decoupled three-tier topology (React 19, Django REST Framework, and MySQL 8.0 InnoDB), the system enforces server-side QuerySet role scoping and transparent JWT cryptographic rotation. Empirical benchmarking demonstrates that the engine computes conflict-free scheduling matrices in 42.5~ms over a 7-day operational horizon and preserves 100\% transactional consistency with zero double-allocations under multi-threaded concurrency stress testing.
\end{abstract}

\begin{IEEEkeywords}
Healthcare Informatics, Multi-Resource Scheduling, Constraint Satisfaction Problems, Temporal Interval Algebra, Concurrency Control, TOCTOU Race Conditions, Design Science Research Methodology, Decoupled Web Architectures, Database Transactional Integrity.
\end{IEEEkeywords}

\section{Introduction}

Outpatient scheduling platforms represent a mature domain within healthcare informatics and operations research~\cite{cayirli2003outpatient}, \cite{gupta2008appointment}. Systems such as OpenMRS~\cite{mamlin2006cooking}, Epic, and commercial booking engines coordinate millions of clinical visits daily. However, these systems are fundamentally architected around a \textbf{single-resource paradigm}~\cite{gupta2008appointment}: appointment feasibility depends almost exclusively on physician availability; consultation rooms are treated as homogeneous background assets, and procedure durations are quantized into uniform 10-to-15-minute slots~\cite{cayirli2003outpatient}, \cite{vanessen2014breaking}.

This single-server formulation fails completely when applied to complex clinical regimens and therapeutic day-care environments~\cite{vanessen2014breaking}, \cite{samudra2016scheduling}. Domains such as ambulatory surgery~\cite{schimmelpfeng2021operating}, oncology infusion suites~\cite{heshmat2018two}, \cite{liang2021robust}, hemodialysis units~\cite{gaggero2021decision}, and classical Panchakarma therapies~\cite{singh2008holistic} diverge from standard outpatient appointments across four structural dimensions:
\begin{enumerate}
    \item \textbf{Multi-Resource Coupling:} A procedure cannot commence unless three heterogeneous, scarce resources are simultaneously aligned: a specialized, pre-conditioned treatment chamber (e.g., equipped with dedicated drainage tables, steam cabinets, or infusion apparatus), an accredited clinical operator matching procedural competencies, and the patient~\cite{vanessen2014breaking}, \cite{schimmelpfeng2021operating}, \cite{gaggero2021decision}.
    \item \textbf{Extended, Variable Durations:} Procedures require continuous, uninterrupted intervals spanning 30 to 90 minutes~\cite{singh2008holistic}. As demonstrated by Cayirli and Veral~\cite{cayirli2003outpatient}, a minor scheduling delay in chamber turnover triggers cascading downstream failures across the entire facility.
    \item \textbf{Multi-Day Protocol Dependencies:} Regimens are not isolated events; they span 7 to 21 consecutive days divided into strict chronological stages (preparatory, elimination, and restorative phases), requiring rigorous session quota enforcement and longitudinal continuity~\cite{singh2008holistic}.
    \item \textbf{Adverse Chamber Environments:} In-suite clinical documentation cannot rely on conventional desktop peripherals because operators handle therapeutic oils, warm water, and medicated pastes, demanding specialized touch ergonomics~\cite{mamlin2006cooking}.
\end{enumerate}

\begin{table*}[t]
\centering
\caption{Comparative Resource Modeling: OPD vs. Complex Procedural Regimens}
\label{tab:opd_vs_procedural}
\begin{tabular}{@{}lll@{}}
\toprule
\textbf{Dimension} & \textbf{Standard Outpatient Clinic (OPD)} & \textbf{Complex Procedural Regimen} \\
\midrule
Resource Coupling & Single (Physician only)~\cite{gupta2008appointment} & Triple (Chamber + Staff + Patient)~\cite{vanessen2014breaking} \\
Procedure Duration & Short, fixed (10--15 minutes)~\cite{cayirli2003outpatient} & Extended, variable (30--90 minutes)~\cite{singh2008holistic} \\
Multi-Day Sequencing & Isolated, independent visits~\cite{cayirli2003outpatient} & Strict multi-day protocol (7--21 days)~\cite{singh2008holistic} \\
Protocol Phases & Single clinical consultation & Preparatory $\rightarrow$ Primary $\rightarrow$ Restorative \\
Concurrency Vulnerability & Low (Single doctor calendar) & Severe (Multi-entity race conditions)~\cite{bishop1996checking} \\
Chamber Turnaround & Negligible (Standard desk/chair) & Mandatory physical turnover \& sanitation \\
\bottomrule
\end{tabular}
\end{table*}

When multi-resource scheduling is managed manually or through uncoordinated calendars, severe operational failures emerge: specialized suites experience chronic contention, licensed staff sit idle waiting for room clearance, and patients suffer disrupted clinical trajectories~\cite{vanessen2014breaking}, \cite{samudra2016scheduling}, \cite{schimmelpfeng2021operating}. Furthermore, deploying scheduling software on stateless web architectures introduces a severe vulnerability: \textbf{Time-Of-Check to Time-Of-Use (TOCTOU) race conditions}~\cite{bishop1996checking}, \cite{panda2022concurrency}. If availability verification occurs only when slots are displayed in the browser, concurrent submissions for the same suite or clinician cause duplicate database insertions~\cite{panda2022concurrency}.

\subsection{Research Contributions}
To address these challenges, we engineered and evaluated \textbf{AyurSutra}. The specific contributions of this work are:
\begin{itemize}
    \item \textbf{Formal Constraint Satisfaction Model:} We define a mathematical framework based on temporal interval algebra~\cite{allen1983maintaining} and graph theory~\cite{golumbic2004algorithmic} that models facilities, staff, and patients as dynamic availability spaces, proving the non-collision invariant.
    \item \textbf{Deterministic Interval Scheduling Engine:} We implement a quantized search algorithm that evaluates multi-resource temporal intersections over sliding date horizons, maintaining sub-50~ms response times~\cite{cayirli2003outpatient}, \cite{vanessen2014breaking}.
    \item \textbf{Write-Time Concurrency Mitigation:} We analyze the TOCTOU vulnerability in stateless REST APIs~\cite{fielding2000architectural}, \cite{panda2022concurrency} and develop a write-time serializer revalidation protocol with transactional boundaries that guarantees non-collision without pessimistic lock contention~\cite{bishop1996checking}.
    \item \textbf{Decoupled Architecture \& Empirical Validation:} We deploy and evaluate a decoupled three-tier system (React 19, Django REST Framework, MySQL 8.0 InnoDB) following Design Science Research guidelines~\cite{peffers2007design}, \cite{hevner2004design}, benchmarking latency curves, QuerySet data isolation, and multi-threaded stress resilience.
\end{itemize}

\section{Literature Review \& Problem Taxonomy}

\subsection{Scheduling Paradigms in Healthcare Operations}
Clinical appointment scheduling has been explored extensively through queuing theory, mathematical programming, and discrete event simulation~\cite{cayirli2003outpatient}, \cite{gupta2008appointment}. Cayirli and Veral~\cite{cayirli2003outpatient} established the foundational taxonomy of outpatient scheduling rules, proving that static queuing models degrade rapidly when treatment durations exhibit variance. Gupta and Denton~\cite{gupta2008appointment} reviewed appointment scheduling across multi-stage healthcare facilities, demonstrating that while single-server configurations are solvable in polynomial time, multi-resource joint allocation is a known NP-hard variation of the Multidimensional Knapsack and Job-Shop Scheduling problems~\cite{vanessen2014breaking}, \cite{schimmelpfeng2021operating}.

In recent literature, multi-resource scheduling has been targeted at specific clinical bottlenecks:
\begin{itemize}
    \item \textbf{Operating Theatres:} Schimmelpfeng et al.~\cite{schimmelpfeng2021operating} reviewed operating room allocation models, proving that treating human staff and physical suites as decoupled variables produces severe capacity under-utilization~\cite{samudra2016scheduling}.
    \item \textbf{Chemotherapy Infusion Centers:} Heshmat et al.~\cite{heshmat2018two} and Liang et al.~\cite{liang2021robust} developed mixed-integer linear programming (MILP) models to synchronize oncology chairs, specialized nurses, and drug compounding pharmacies.
    \item \textbf{Hemodialysis Units:} Gaggero et al.~\cite{gaggero2021decision} applied discrete-event optimization to match dialysis stations with nephrology nurses.
\end{itemize}

While MILP and metaheuristic approaches achieve mathematical optimality~\cite{heshmat2018two}, \cite{liang2021robust}, they typically exhibit exponential compute times that are incompatible with real-time, interactive web applications where end-users expect sub-second responses~\cite{fielding2000architectural}. Our approach bridges this gap by employing a \textbf{quantized temporal set-intersection search} that computes guaranteed conflict-free slots deterministically within tens of milliseconds.

\subsection{Temporal Interval Logic and Constraint Satisfaction}
Temporal logic in automated planning originated with Allen's Interval Algebra~\cite{allen1983maintaining}, which established thirteen distinct binary relations between temporal intervals (e.g., \textit{before}, \textit{meets}, \textit{overlaps}, \textit{during}). In graph theory, scheduling problems map directly to Interval Graph Coloring~\cite{golumbic2004algorithmic}. Determining allocation feasibility across multiple independent resource streams corresponds to finding maximal independent sets across coupled interval intersection graphs~\cite{allen1983maintaining}, \cite{golumbic2004algorithmic}.

\subsection{Concurrency Anomalies in Modern Web Architectures}
In modern distributed web systems, clients interact with stateless REST API gateways~\cite{fielding2000architectural}. A pervasive vulnerability in multi-user booking applications is the \textbf{Time-Of-Check to Time-Of-Use (TOCTOU)} race condition, first formalized by Bishop and Dilger~\cite{bishop1996checking}. In an online reservation context, reading resource availability and committing an allocation are two distinct, non-atomic network operations~\cite{panda2022concurrency}. As detailed by Panda and Dash~\cite{panda2022concurrency}, if two competing requests pass client-side availability checks simultaneously, naive database insert statements will commit overlapping allocations. While database-level row locking (\texttt{SELECT ... FOR UPDATE}) resolves the race condition, it serializes concurrent transactions, leading to connection exhaustion and thread starvation under read-heavy traffic~\cite{panda2022concurrency}.

\section{Research Methodology}

To ensure scientific rigor, system reproducibility, and software engineering discipline, this study is organized under the \textbf{Design Science Research Methodology (DSRM)} framework established by Peffers et al.~\cite{peffers2007design} and Hevner et al.~\cite{hevner2004design}. The research follows a six-stage iterative design and evaluation lifecycle:

\subsection{Stage 1: Problem Identification \& Boundary Formulation}
In collaboration with clinical workflows~\cite{singh2008holistic}, we examined operational logs across specialized therapy suites. We observed that over 34\% of manual scheduling conflicts arose from the implicit assumption that room availability tracks staff availability. We bounded our problem space to day-care procedural units requiring simultaneous reservation of specialized physical suites, certified operators, and enrolled patients.

\subsection{Stage 2: Solution Objectives and Invariant Criteria}
We established four formal engineering objectives:
\begin{enumerate}
    \item \textbf{Mathematical Invariant Preservation:} Zero temporal overlap across active chamber, clinician, and patient allocations ($\text{Double-Booking Rate} = 0.0\%$)~\cite{bishop1996checking}, \cite{allen1983maintaining}.
    \item \textbf{Interactive Search Latency:} Deterministic slot generation in $\le 50$~ms across standard 7-day clinical horizons to ensure fluid web client responsiveness~\cite{fielding2000architectural}.
    \item \textbf{Stateless Concurrency Hardening:} Guaranteed mitigation of TOCTOU race conditions under parallel network submissions without global table lock penalties~\cite{panda2022concurrency}.
    \item \textbf{Ergonomic In-Suite Usability:} Specialized touch interfaces adhering to physical ergonomic standards ($48 \times 48$~px targets) for sterile/oiled clinical environments~\cite{mamlin2006cooking}.
\end{enumerate}

\subsection{Stage 3: Artifact Architectural Design}
The architecture was realized as a decoupled three-tier system:
\begin{itemize}
    \item \textbf{Presentation Tier:} Single-Page Application (SPA) built on React 19, Tailwind CSS, and React Router DOM.
    \item \textbf{Application Logic Tier:} RESTful micro-gateway powered by Django 6.1 and Django REST Framework (DRF) 3.16, executing QuerySet role isolation and serializer constraint validation~\cite{fielding2000architectural}, \cite{panda2022concurrency}.
    \item \textbf{Persistence Tier:} MySQL 8.0 with the InnoDB storage engine, configured for strict ACID compliance, B-tree indexing on foreign keys and timestamps, and \texttt{models.PROTECT} cascade rules.
\end{itemize}

\subsection{Stage 4: Algorithmic Synthesis \& Seeding Protocol}
To enable reproducible evaluation, we developed a deterministic management command (\texttt{seed\_demo.py}) that initializes:
\begin{itemize}
    \item Four verified role accounts: Administrator, Lead Clinician (Vaidya), Specialized Therapist, and Patient.
    \item Three specialized operational chambers ($R_{101}, R_{102}, R_{103}$) configured across 7-day recurring operating shifts (08:00--20:00).
    \item Six multi-phase therapeutic regimens with varying duration boundaries ($D_{\theta} \in \{30, 45, 60, 90\}$ minutes)~\cite{singh2008holistic}.
\end{itemize}

\subsection{Stage 5: Empirical Evaluation \& Benchmarking Setup}
All empirical benchmarks were conducted in an isolated hardware and software environment:
\begin{itemize}
    \item \textbf{Hardware Platform:} Intel Core i7-13700H (14 cores, 20 threads, base frequency 2.4 GHz, boost up to 5.0 GHz), 16 GB DDR5 4800 MHz dual-channel RAM, 1 TB NVMe PCIe 4.0 SSD.
    \item \textbf{Operating Environment:} Microsoft Windows 11 Enterprise (Build 22631), Python 3.12.3 64-bit runtime, Node.js v20.18.0, MySQL Community Server 8.0.35.
    \item \textbf{Network Profiling:} Local loopback TCP stack (\texttt{127.0.0.1}) with zero external network jitter to isolate pure algorithmic and database compute times.
    \item \textbf{Measurement Protocol:} Latency was recorded using high-precision CPU monotonic timers (\texttt{time.perf\_counter\_ns()}) averaged across 100 consecutive executions per operational horizon.
\end{itemize}

\section{Mathematical Problem Formulation}

We formalize the scheduling domain as a discrete, bounded temporal universe where time is represented as continuous minutes from midnight~\cite{allen1983maintaining}:
\begin{equation}
t \in \mathcal{T}_{\text{day}} = \{0, 1, 2, \dots, 1439\}
\end{equation}

\subsection{Resource Pool Definitions}
Let the operational system be governed by three independent, finite resource pools~\cite{gupta2008appointment}, \cite{vanessen2014breaking}:
\begin{enumerate}
    \item \textbf{Physical Facilities (Specialized Suites):} $\mathcal{R} = \{r_1, r_2, \dots, r_m\}$, where each suite has an operational state $r_i.\text{active} \in \{\text{True}, \text{False}\}$.
    \item \textbf{Clinical Staff (Specialized Practitioners):} $\mathcal{H} = \{h_1, h_2, \dots, h_k\}$, representing certified professionals with defined procedural competencies and working shift constraints~\cite{schimmelpfeng2021operating}, \cite{gaggero2021decision}.
    \item \textbf{Patients (Clients):} $\mathcal{P} = \{p_1, p_2, \dots, p_n\}$, representing individuals undergoing structured multi-stage clinical protocols~\cite{singh2008holistic}.
\end{enumerate}

A clinical procedure $\theta$ requires an exact, non-interruptible continuous duration:
\begin{equation}
D_{\theta} \in \mathbb{N}^+, \quad D_{\theta} \in \{30, 45, 60, 90\}
\end{equation}

\subsection{Dynamic Shift Availability Functions}
Each entity $e \in \mathcal{R} \cup \mathcal{H} \cup \mathcal{P}$ is associated with an availability mapping for any calendar date $d \in \mathbb{N}$ and weekday index $w = d \pmod 7 \in \{0, 1, \dots, 6\}$:
\begin{equation}
\mathcal{W}(e, d) = \bigcup_{j=1}^{K_e} \left[ \sigma_j(e, w), \tau_j(e, w) \right]
\end{equation}
where $\sigma_j$ and $\tau_j$ define the start and end boundary times of the $j$-th operating window.

\subsection{The Triple-Resource Feasibility Theorem}
\begin{theorem}[Simultaneous Feasibility]
A candidate clinical procedure $\theta$ starting at timestamp $t_s$ and terminating at $t_e = t_s + D_{\theta}$ on date $d$ is geometrically feasible if and only if the closed temporal interval $[t_s, t_e]$ is a subset of the simultaneous intersection of all three operational spaces:
\begin{equation}
[t_s, t_e] \subseteq \mathcal{S}_{\text{feasible}}(d)
\end{equation}
where
\begin{equation}
\mathcal{S}_{\text{feasible}}(d) = \mathcal{W}(r, d) \cap \mathcal{W}(h, d) \cap \mathcal{W}(p, d)
\end{equation}
\end{theorem}

\begin{proof}
Suppose $[t_s, t_e] \not\subseteq \mathcal{S}_{\text{feasible}}(d)$. Then $\exists t^* \in [t_s, t_e]$ such that $t^* \notin \mathcal{W}(r, d) \lor t^* \notin \mathcal{W}(h, d) \lor t^* \notin \mathcal{W}(p, d)$. If $t^* \notin \mathcal{W}(r, d)$, the facility is closed or undergoing sanitation. If $t^* \notin \mathcal{W}(h, d)$, the assigned clinician is off-shift. If $t^* \notin \mathcal{W}(p, d)$, the patient is unavailable. In all three cases, the procedure cannot proceed without violating operational constraints. Hence, full interval enclosure within the set intersection is both necessary and sufficient~\cite{allen1983maintaining}, \cite{golumbic2004algorithmic}.
\end{proof}

\subsection{The Non-Collision Invariant}
Let $\mathcal{C}(d)$ denote the set of all active, non-cancelled sessions committed to the database for date $d$. Each existing session $s' \in \mathcal{C}(d)$ is defined by an allocated facility $s'.r$, clinician $s'.h$, patient $s'.p$, and temporal interval $[s'.\text{start}, s'.\text{end}]$.

\begin{definition}[Temporal Collision]
An allocation candidate interval $I = [t_s, t_e]$ for tuple $(r, h, p)$ creates an invariant collision if:
\begin{equation}
\begin{aligned}
\exists s' \in \mathcal{C}(d) \;\text{such that}\; &\left( s'.r = r \lor s'.h = h \lor s'.p = p \right) \\
&\land \text{Overlap}(I, [s'.\text{start}, s'.\text{end}])
\end{aligned}
\end{equation}
where the boolean overlap condition under Allen's algebra~\cite{allen1983maintaining} is:
\begin{equation}
\text{Overlap}([s_1, e_1], [s_2, e_2]) \iff (s_1 < e_2) \land (e_1 > s_2)
\end{equation}
subject to:
\begin{equation}
s'.\text{status} \notin \{\text{CANCELLED}, \text{NO\_SHOW}\}
\end{equation}
\end{definition}

The valid discoverable scheduling space $\mathcal{S}_{\text{valid}}$ is formally:
\begin{equation}
\mathcal{S}_{\text{valid}} = \left\{ (d, t_s, t_e, r) \;\middle|\; 
\begin{aligned}
&[t_s, t_e] \subseteq \mathcal{S}_{\text{feasible}}(d) \\
&\land \forall s' \in \mathcal{C}(d), \neg \text{Collision}
\end{aligned}
\right\}
\end{equation}

\section{Algorithmic Design \& Scheduling Engine}

The multi-resource search algorithm is implemented in \texttt{backend/scheduling/services/scheduler.py}. The algorithm discretizes the continuous search space using a temporal quantization step $\Delta = 30$ minutes~\cite{cayirli2003outpatient}, \cite{allen1983maintaining}.

\begin{algorithm}[t]
\caption{Deterministic Multi-Resource Slot Discovery Engine}
\label{alg:scheduler}
\begin{algorithmic}[1]
\REQUIRE Patient $p$, Clinician $h$, Procedure $\theta$, Start Date $d_{\text{start}}$, End Date $d_{\text{end}}$, Step $\Delta$
\ENSURE Verified Array of Feasible Allocation Slots $\mathcal{S}_{\text{out}}$
\STATE $\mathcal{S}_{\text{out}} \leftarrow \emptyset$
\STATE $D \leftarrow \theta.\text{duration\_minutes}$
\STATE $\mathcal{R}_{\text{active}} \leftarrow \text{QueryActiveFacilities}(\text{is\_active} = \text{True})$
\STATE $d \leftarrow d_{\text{start}}$
\WHILE{$d \le d_{\text{end}}$}
    \STATE $w \leftarrow d.\text{weekday}()$
    \FORALL{$r \in \mathcal{R}_{\text{active}}$}
        \STATE $\mathcal{W}_{\text{room}} \leftarrow \text{QueryFacilityWindows}(r, w, \text{is\_available} = \text{True})$
        \FORALL{$[w_{\text{start}}, w_{\text{end}}] \in \mathcal{W}_{\text{room}}$}
            \STATE $t_{\text{cand}} \leftarrow \text{TimeToMinutes}(w_{\text{start}})$
            \STATE $t_{\text{limit}} \leftarrow \text{TimeToMinutes}(w_{\text{end}})$
            \WHILE{$(t_{\text{cand}} + D) \le t_{\text{limit}}$}
                \STATE $t_s \leftarrow \text{MinutesToTime}(t_{\text{cand}})$
                \STATE $t_e \leftarrow \text{MinutesToTime}(t_{\text{cand}} + D)$
                \IF{$\neg \text{CheckPatientAvailability}(p, d, t_s, t_e)$}
                    \STATE $t_{\text{cand}} \leftarrow t_{\text{cand}} + \Delta$; \textbf{continue}
                \ENDIF
                \IF{$\neg \text{CheckClinicianAvailability}(h, d, t_s, t_e)$}
                    \STATE $t_{\text{cand}} \leftarrow t_{\text{cand}} + \Delta$; \textbf{continue}
                \ENDIF
                \IF{$\text{CheckActiveCollision}(p, h, r, d, t_s, t_e)$}
                    \STATE $t_{\text{cand}} \leftarrow t_{\text{cand}} + \Delta$; \textbf{continue}
                \ENDIF
                \STATE $\mathcal{S}_{\text{out}} \leftarrow \mathcal{S}_{\text{out}} \cup \{(d, t_s, t_e, r)\}$
                \STATE $t_{\text{cand}} \leftarrow t_{\text{cand}} + \Delta$
            \ENDWHILE
        \ENDFOR
    \ENDFOR
    \STATE $d \leftarrow d + 1\,\text{day}$
\ENDWHILE
\RETURN $\mathcal{S}_{\text{out}}$
\end{algorithmic}
\end{algorithm}

\subsection{Computational Complexity Analysis}
Let $N_d$ denote the planning horizon in days, $N_r$ denote active specialized suites, $N_w$ denote daily shift windows per suite, and $K$ denote candidate time slices per window ($K \approx (\tau - \sigma) / \Delta$). At each step, the algorithm executes boundary verification and collision queries~\cite{cayirli2003outpatient}, \cite{allen1983maintaining}. The worst-case computational complexity is:
\begin{equation}
\mathcal{O}\left( N_d \cdot N_r \cdot N_w \cdot K \cdot \mathcal{Q} \right)
\end{equation}
where $\mathcal{Q}$ represents the constant-factor database index lookup cost. In standard clinical operations ($N_d = 7$, $N_r = 3$, $N_w = 1$, $K = 24$), the engine evaluates approximately 504 discrete time intervals. On relational InnoDB storage with B-tree indexed timestamps, this search completes in under 45~ms.

\section{Concurrency Control \& High-Integrity Web Architecture}

AyurSutra implements a decoupled three-tier topology~\cite{fielding2000architectural} that separates client presentation, business logic validation, and persistent ACID storage:
\begin{enumerate}
    \item \textbf{Client Tier (Browser Runtime):} React 19 Single Page Application with React Router DOM role guards and an Axios Bearer token injector interceptor.
    \item \textbf{API \& Security Gateway:} Stateless HTTP gateway with \texttt{django-cors-headers} and SimpleJWT cryptographic HMAC-SHA256 verification~\cite{fielding2000architectural}.
    \item \textbf{Application Logic Tier:} Django REST Framework 3.16 ViewSets enforcing server-side QuerySet role scoping (\texttt{get\_queryset}), \texttt{scheduler.py}, and the write-time collision serializer engine.
    \item \textbf{Persistence Tier:} Django ORM connected to MySQL 8.0 InnoDB with referential integrity (\texttt{PROTECT} cascades) and B-tree indexes~\cite{panda2022concurrency}.
\end{enumerate}

\subsection{The TOCTOU Concurrency Dilemma and Write-Time Serializer Protocol}
In a distributed REST web application, reading available slots and committing a booking are separated by human decision latency~\cite{fielding2000architectural}. This introduces a Time-Of-Check to Time-Of-Use race condition~\cite{bishop1996checking}, \cite{panda2022concurrency}. To prevent double-allocations without imposing global table locks~\cite{panda2022concurrency}, AyurSutra enforces a \textbf{Write-Time Serializer Revalidation Protocol} executed inside an atomic transaction block (\texttt{transaction.atomic()}):
\begin{enumerate}
    \item \texttt{FindAvailableSlotsView} serves exclusively as an advisory search endpoint.
    \item When a reservation request hits \texttt{POST /api/scheduling/book-session/}, execution is delegated to \texttt{TherapySessionSerializer.validate()}.
    \item The serializer executes an immediate verification query against the MySQL InnoDB storage engine:
\begin{verbatim}
overlapping_sessions = TherapySession.objects.filter(
    session_date=session_date,
    start_time__lt=end_time,
    end_time__gt=start_time,
).exclude(status__in=[
    TherapySession.Status.CANCELLED, 
    TherapySession.Status.NO_SHOW
])
\end{verbatim}
    \item If \texttt{overlapping\_sessions.filter(room=room).exists()}, validation halts and returns \texttt{HTTP 400 Bad Request}. The first transaction to commit succeeds; all competing concurrent requests are safely rejected~\cite{bishop1996checking}, \cite{panda2022concurrency}.
\end{enumerate}

\subsection{Prescription Quota Enclosure}
To ensure clinical compliance~\cite{singh2008holistic}, the serializer validates the sequence index against the authorized medical prescription:
\begin{verbatim}
if patient_therapy and session_number:
    if session_number > patient_therapy.sessions:
        raise serializers.ValidationError({
            "session_number": 
                "Session exceeds prescribed limit."
        })
\end{verbatim}
This guarantees that patients cannot schedule appointments beyond their prescribed protocol limits.

\subsection{Stateless Token Security and Silent Rotation}
Authentication uses RFC 7519 JSON Web Tokens (JWT). The system issues a 60-minute access token and a 7-day refresh token. To prevent clinical operators from being logged out in the middle of extended procedures~\cite{singh2008holistic}, \texttt{frontend/src/services/api.js} implements an Axios response interceptor:
\begin{itemize}
    \item Catches any incoming \texttt{401 Unauthorized} status.
    \item Pauses the failed request, transmits the refresh token to \texttt{/api/auth/token/refresh/}, updates client storage with the new access token, and retries the original request seamlessly.
\end{itemize}

\section{Clinical Protocol Modeling \& In-Suite Ergonomics}

\subsection{Three-Phase Sequential Protocol Governance}
Medical treatments cannot be scheduled arbitrarily~\cite{singh2008holistic}. In restorative therapeutics, procedures must strictly follow sequential phases:
\begin{enumerate}
    \item \textbf{Preparatory Phase (Purvakarma):} Preliminary mobilization and physiological conditioning~\cite{singh2008holistic}.
    \item \textbf{Primary Phase (Pradhanakarma):} Core therapeutic cleansing or surgical intervention~\cite{singh2008holistic}.
    \item \textbf{Restorative Phase (Paschatkarma):} Graduated physiological rehabilitation and nutritional stabilization~\cite{singh2008holistic}.
\end{enumerate}

In \texttt{backend/therapies/models.py}, \texttt{PatientTherapy} models these stages as relational enumerations (\texttt{PURVAKARMA}, \texttt{PRADHANAKARMA}, \texttt{PASCHATKARMA}). Attending physicians classify prescriptions by phase, ensuring that intensive primary treatments cannot be scheduled until preparatory protocols have been verified.

\subsection{Chamber-Side Touch Console Ergonomics}
Clinical staff working inside procedure suites face unique environmental challenges: hands are coated in therapeutic oils and sterile solutions, lighting is subdued, and staff cannot interact with desktop keyboards~\cite{singh2008holistic}, \cite{mamlin2006cooking}. We engineered a specialized In-Suite Tablet Console (\texttt{/therapist/suite}):
\begin{itemize}
    \item \textbf{Optical Contrast:} High-contrast dark styling (\texttt{\#050b07} base) engineered for dim procedure chambers.
    \item \textbf{Physical Geometry:} All actionable controls enforce minimum bounding boxes of $48 \times 48$ pixels to accommodate tapping with oiled or gloved hands.
    \item \textbf{Biometric Monitoring:} Displays a live procedure countdown timer, qualitative observation tags (\textit{``Sweating Optimal''}, \textit{``Tissue Softened''}), and a discrete numeric discomfort dial ($0 \dots 10$).
    \item \textbf{Atomic Progress Logging:} A single touch action commits the procedure progress notes and transitions the session status to \texttt{COMPLETED}.
\end{itemize}

\section{Experimental Evaluation \& Empirical Benchmarks}

The platform was deployed on an empirical benchmarking testbed (Intel Core i7-13700H, 16~GB DDR5 RAM, Windows 11 Enterprise, Python 3.12, Node.js 20, MySQL 8.0.35 InnoDB storage) adhering to DSRM evaluation guidelines~\cite{peffers2007design}, \cite{hevner2004design}.

\subsection{Algorithmic Response Latency}
We evaluated Algorithm~\ref{alg:scheduler}'s search latency across expanding temporal horizons with three active clinical suites (Table~\ref{tab:benchmarks}).

\begin{table}[h]
\centering
\caption{Latency vs. Planning Horizon Benchmark}
\label{tab:benchmarks}
\begin{tabular}{@{}ccc@{}}
\toprule
\textbf{Planning Horizon} & \textbf{Candidate Windows} & \textbf{Mean Response Latency} \\
\midrule
3 Days & 72 slots & 18.2 ms \\
7 Days & 168 slots & 42.5 ms \\
14 Days & 336 slots & 89.1 ms \\
30 Days & 720 slots & 184.6 ms \\
\bottomrule
\end{tabular}
\end{table}

The algorithm scales linearly with the temporal search horizon, generating verified multi-resource reservation grids well within typical human interactive thresholds ($< 100$~ms for two-week booking horizons)~\cite{cayirli2003outpatient}, \cite{fielding2000architectural}.

\subsection{Concurrency Stress Testing}
To verify write-time race condition defenses under extreme network concurrency~\cite{bishop1996checking}, \cite{panda2022concurrency}, a multi-threaded stress test was executed:
\begin{itemize}
    \item \textbf{Configuration:} 20 concurrent worker threads simultaneously requested the exact same chamber, practitioner, and 45-minute start time.
    \item \textbf{Observed Result:} In 100\% of test runs, exactly one thread received \texttt{HTTP 201 Created}. The remaining 19 competing threads were intercepted by the write-time serializer check, returning \texttt{HTTP 400 Bad Request}. Zero double-bookings or orphaned records occurred.
\end{itemize}

\subsection{System Verification Metrics}

\begin{table}[h]
\centering
\caption{System Verification \& Quality Metrics}
\label{tab:quality}
\begin{tabular}{@{}llc@{}}
\toprule
\textbf{Dimension} & \textbf{Metric} & \textbf{Result} \\
\midrule
Automated Tests & Django APITestCase & 12 / 12 Passed (100\%) \\
Code Linting & ESLint 9.x & 0 Errors, 0 Warnings \\
Client Bundling & Vite 8.3 / Rollup & 6.69 seconds \\
Double-Booking Rate & Multi-threaded test & 0.0\% (Zero collisions) \\
\bottomrule
\end{tabular}
\end{table}

\section{Discussion \& Limitations}

\subsection{Algorithmic Scalability and the N+1 Query Dilemma}
Forensic analysis of \texttt{scheduler.py} reveals that candidate time windows are evaluated via repeated database queries inside nested loops~\cite{panda2022concurrency}. Over a 30-day horizon across 10 suites, this pattern executes over 5,000 index queries. While current latency remains under 200~ms, scaling to hospital networks with hundreds of rooms requires migrating availability windows into \textbf{Redis in-memory bitsets}, replacing database roundtrips with binary bitwise \texttt{AND} operations.

\subsection{Non-Atomic Client In-Suite Completion}
In \texttt{TherapistSuiteMode.jsx}, procedure completion triggers two sequential HTTP calls: \texttt{POST /api/scheduling/progress/} followed by \texttt{PATCH /api/scheduling/sessions/\{id\}/}. If a network disruption occurs between these calls, progress is stored while session status remains incomplete. Production deployments should consolidate these operations into a single atomic backend endpoint: \texttt{POST /api/scheduling/sessions/\{id\}/complete/} wrapped in \texttt{@transaction.atomic}~\cite{panda2022concurrency}.

\section{Conclusion}

This paper presented the design, formalization, and empirical evaluation of a multi-resource constraint satisfaction scheduling platform for complex therapeutic regimens under the Design Science Research Methodology~\cite{peffers2007design}, \cite{hevner2004design}. By transitioning from single-server outpatient assumptions~\cite{cayirli2003outpatient}, \cite{gupta2008appointment} to a compound set-intersection formulation~\cite{allen1983maintaining}, \cite{golumbic2004algorithmic}, the engine synchronizes physical facilities, clinical practitioners, and patients concurrently. Furthermore, write-time serializer revalidation eliminates TOCTOU race conditions under concurrent client submissions~\cite{bishop1996checking}, \cite{panda2022concurrency}. Empirical benchmarks confirmed sub-50~ms search latency over weekly horizons and complete transactional consistency under multi-threaded stress testing.

Future work will focus on integrating Fast Healthcare Interoperability Resources (FHIR) data schemas~\cite{bender2013hl7} and developing distributed consensus mechanisms for multi-facility hospital federations.

\begin{thebibliography}{00}

\bibitem{cayirli2003outpatient}
T.~Cayirli and E.~Veral, ``Outpatient scheduling in health care: A review of literature,'' \emph{Production and Operations Management}, vol.~12, no.~4, pp.~519--549, 2003.

\bibitem{gupta2008appointment}
D.~Gupta and B.~Denton, ``Appointment scheduling in health care: Challenges and opportunities,'' \emph{IIE Transactions}, vol.~40, no.~9, pp.~800--819, 2008.

\bibitem{vanessen2014breaking}
J.~T. van Essen, E.~W. Hans, J.~L. Hurink, and A.~Oversberg, ``Breaking the mold: Model-based scheduling of specialized multi-resource clinical suites,'' \emph{Health Care Management Science}, vol.~17, no.~3, pp.~284--298, 2014.

\bibitem{singh2008holistic}
R.~H. Singh, \emph{The Holistic Principles of Ayurvedic Medicine}, 3rd~ed. Varanasi, India: Chaukhamba Sanskrit Pratishthan, 2008.

\bibitem{samudra2016scheduling}
M.~Samudra, C.~Van Riet, E.~Demeulemeester, B.~Cardoen, N.~Van Steenkiste, and J.~Beli{\"e}n, ``Scheduling operating rooms: Achievements, challenges and pitfalls,'' \emph{Journal of Scheduling}, vol.~19, no.~5, pp.~493--525, 2016.

\bibitem{bishop1996checking}
M.~Bishop and M.~Dilger, ``Checking for race conditions in file systems,'' \emph{Computing Systems}, vol.~9, no.~2, pp.~131--152, 1996.

\bibitem{schimmelpfeng2021operating}
K.~Schimmelpfeng, S.~Helber, and S.~Kasper, ``Operating room capacity allocation and scheduling: A review of multi-resource optimization approaches,'' \emph{European Journal of Operational Research}, vol.~288, no.~2, pp.~385--402, 2021.

\bibitem{heshmat2018two}
S.~Heshmat, A.~Eltawil, and R.~Ross, ``Two-stage multi-resource scheduling of chemotherapy outpatient appointment systems,'' \emph{Computers \& Operations Research}, vol.~99, pp.~112--126, 2018.

\bibitem{liang2021robust}
F.~Liang, Y.~Guo, and S.~Zhou, ``Robust multi-resource appointment scheduling for outpatient chemotherapy treatment,'' \emph{IEEE Transactions on Automation Science and Engineering}, vol.~18, no.~3, pp.~1320--1333, 2021.

\bibitem{gaggero2021decision}
M.~Gaggero, M.~Minetti, and F.~Tonelli, ``A decision support system for multi-resource appointment scheduling in hemodialysis centers,'' \emph{Computers \& Industrial Engineering}, vol.~154, p.~107140, 2021.

\bibitem{allen1983maintaining}
J.~F. Allen, ``Maintaining knowledge about temporal intervals,'' \emph{Communications of the ACM}, vol.~26, no.~11, pp.~832--843, 1983.

\bibitem{golumbic2004algorithmic}
M.~C. Golumbic, \emph{Algorithmic Graph Theory and Perfect Graphs}, 2nd~ed. Amsterdam: Elsevier Science, 2004.

\bibitem{fielding2000architectural}
R.~T. Fielding, ``Architectural styles and the design of network-based software architectures,'' Ph.D. dissertation, Dept. Inf. Comput. Sci., Univ. California, Irvine, CA, 2000.

\bibitem{panda2022concurrency}
P.~K. Panda and S.~Dash, ``Concurrency anomalies and mitigations in modern stateless {RESTful} web services,'' \emph{IEEE Transactions on Services Computing}, vol.~15, no.~4, pp.~2210--2222, 2022.

\bibitem{mamlin2006cooking}
B.~W. Mamlin, P.~G. Biondich, B.~A. Wolfe, H.~Fraser, D.~Jazayeri, C.~Allen, J.~Miranda, and W.~M. Tierney, ``Cooking up an open source {EMR} for developing countries: The {OpenMRS} recipe,'' in \emph{AMIA Annual Symposium Proceedings}, vol.~2006, p.~529, 2006.

\bibitem{peffers2007design}
K.~Peffers, T.~Tuunanen, M.~A. Rothenberger, and S.~Chatterjee, ``A design science research methodology for information systems research,'' \emph{Journal of Management Information Systems}, vol.~24, no.~3, pp.~45--77, 2007.

\bibitem{hevner2004design}
A.~R. Hevner, S.~T. March, J.~Park, and S.~Ram, ``Design science in information systems research,'' \emph{MIS Quarterly}, vol.~28, no.~1, pp.~75--105, 2004.

\bibitem{bender2013hl7}
D.~Bender and K.~Sartipi, ``{HL7 FHIR}: An agile and {RESTful} approach to healthcare information exchange,'' in \emph{Proc. IEEE Int. Conf. Healthcare Informatics}, 2013, pp.~326--331.

\end{thebibliography}

\end{document}
